\documentclass[12pt,a4paper]{book}{}
%\documentclass[titlepage,12pt,a4paper]{book}

% substituir linha seguinte por 
% \usepackage[english]{babel} 
% se o relatório for escrito na língua inglesa.
\usepackage[portuguese]{babel}

% \usepackage[utf8]{inputenc}
\usepackage[T1]{fontenc}

\usepackage{makeidx}
\usepackage{xspace}
\usepackage{graphicx,color,times}
\usepackage{fancyhdr}
% \usepackage{pxfonts}
% \usepackage{times}
% \usepackage{mathptm}
% \usepackage{amssymb}
% \usepackage{amsfonts}

\usepackage{amsmath}
\usepackage{latexsym}
\usepackage[printonlyused]{acronym}
\usepackage{float}
\usepackage{listings}
\usepackage{tocbibind}
\usepackage{natbib}
\usepackage{hyperref}

% \usepackage{glossaries}
% \makeglossaries

% \renewcommand{\ttdefault}{phv}

\pagestyle{fancy}
\renewcommand{\chaptermark}[1]{\markboth{#1}{}}
\renewcommand{\sectionmark}[1]{\markright{\thesection\ #1}}
\fancyhf{} \fancyhead[LE,RO]{\bfseries\thepage}
\fancyhead[LO]{\bfseries\rightmark}
\fancyhead[RE]{\bfseries\leftmark}
\renewcommand{\headrulewidth}{0.5pt}
\renewcommand{\footrulewidth}{0pt}
\setlength{\headheight}{15.6pt}
\setlength{\marginparsep}{0cm}
\setlength{\marginparwidth}{0cm}
\setlength{\marginparpush}{0cm}
\addtolength{\hoffset}{-1.0cm}
\addtolength{\oddsidemargin}{\evensidemargin}
\addtolength{\oddsidemargin}{0.5cm}
\addtolength{\evensidemargin}{-0.5cm}


\usepackage{fix-cm}
\usepackage{fourier}
\usepackage[scaled=.92]{helvet}
% \definecolor{ChapGrey}{rgb}{0.6,0.6,0.6}
% \newcommand{\LargeFont}{
%   \usefont{\encodingdefault}{\rmdefault}{b}{n}
%   \fontsize{60}{80}\selectfont\color{ChapGrey}
%   }
% \makeatletter
% C-c\makechapterstyle{GreyNum}{
%   \renewcommand{\chapnamefont}{\large\sffamily\bfseries\itshape}
%   \renewcommand{\chapnumfont}{\LargeFont}
%   \renewcommand{\chaptitlefont}{\Huge\sffamily\bfseries\itshape}
%   \setlength{\beforechapskip}{0pt}
%   \setlength{\midchapskip}{40pt}
%   \setlength{\afterchapskip}{60pt}
%   \renewcommand\chapterheadstart{\vspace*{\beforechapskip}}
%   \renewcommand\printchaptername{
%   \begin{tabular}{@{}c@{}}
%     \chapnamefont \@chapapp\\}
%     \renewcommand\chapternamenum{\noalign{\vskip 2ex}}
%     \renewcommand\printchapternum{\chapnumfont\thechapter\par}
%     \renewcommand\afterchapternum{
%   \end{tabular}
%   \par\nobreak\vskip\midchapskip}
%   \renewcommand\printchapternonum{}
%   \renewcommand\printchaptertitle[1]{
%   {\chaptitlefont{##1}\par}}
%   \renewcommand\afterchaptertitle{\par\nobreak\vskip \afterchapskip}
% }
% \makeatother
% \chapterstyle{GreyNum}

% \setcounter{tocdepth}{3}
% \setsecnumdepth{subsubsection}

\renewcommand{\ttdefault}{lmtt}


% NEW COLORS
\definecolor{dark}{gray}{0.25}
\definecolor{lgray}{gray}{0.9}
\definecolor{dkblue}{rgb}{0,0.13,0.4}
\definecolor{dkgreen}{rgb}{0,0.6,0}
\definecolor{gray}{rgb}{0.5,0.5,0.5}
\definecolor{mauve}{rgb}{0.58,0,0.82}

\lstset{ %
  language=C,                    basicstyle=\footnotesize,
  numbers=none,                  numberstyle=\tiny\color{gray}, 
  stepnumber=1,                  numbersep=5pt,
  backgroundcolor=\color{white}, showspaces=false,
  showstringspaces=false,        showtabs=false,
  frame=single,                  rulecolor=\color{black},
  tabsize=2,                     captionpos=b,
  breaklines=true,               breakatwhitespace=false,
  title=\lstname,                keywordstyle=\color{blue},
  commentstyle=\color{dkgreen},  stringstyle=\color{mauve},
  escapeinside={\%*}{*)},        morekeywords={*},
  belowskip=0cm
}

\renewcommand{\lstlistingname}{Excerto de Código}
\renewcommand{\lstlistlistingname}{Lista de Excertos de Código}

\renewcommand{\today}{\day \ifcase \month \or Janeiro\or Fevereiro\or Março\or %
Abril\or Maio\or Junho\or Julho\or Agosto\or Setembro\or Outubro\or Novembro\or %
Dezembro\fi de \number \year} 



\begin{document}


\thispagestyle{empty}
\setcounter{page}{-1}

\begin{center}
\begin{Huge}
\textbf{Universidade da Beira Interior}
\end{Huge}
\end{center}

\begin{center}
\begin{Huge}
Departamento de Informática
\end{Huge}
\end{center}

\vspace{0,07cm}
\begin{figure}[!htb]
\centering
\includegraphics[width=191pt]{ubi-fe-di.png}
\end{figure}

\vspace{0.5cm}
\begin{center}
\begin{Large}
\textbf{QUEM GOVERNA A INTERNET}
\end{Large}
\end{center}


\vspace{0.5cm}
\begin{center}
\begin{normalsize}
\begin{large}
Elaborado por:
\end{large}
\end{normalsize}
\end{center}

\vspace{0.2cm}
\begin{center}
\begin{large}
\textbf{Leonardo Ferreira [52981]}
\end{large}
\end{center}

\vspace{0,5cm}
\begin{center}
\begin{normalsize}
\begin{large}
Orientador:
\end{large}
\end{normalsize}
\end{center}

\vspace{0.2cm}
\begin{center}
\begin{large}
\textbf{Professor Doutor Bruno Miguel Silva}
\end{large}
\end{center}



\vspace{0.5cm}
\begin{center}
\begin{normalsize}
\today
\end{normalsize}
\end{center}

\let\cleardoublepage\clearpage
\chapter*{Agradecimentos}
\pagenumbering{arabic}
\label{chap:ack}

\paragraph{} A conclusão deste trabalho
%\frontmatte
\tableofcontents

%\listoffigures

% #   ATENÇÃO
% Se não existirem tabelas, comentar as duas linhas seguintes
%\clearpage{\thispagestyle{empty}\cleardoublepage}
%\listoftables

% #   ATENÇÃO
% Se existirem trechos de código, descomentar as seguintes linhas
% \clearpage{\thispagestyle{empty}\cleardoublepage}
% \lstlistoflistings

\chapter*{Acrónimos}
\addcontentsline{toc}{chapter}{Acrónimos}
\label{chap:acro}

\begin{acronym}[UBI]

  \acro{UBI}{Universidade da Beira Interior}

\end{acronym}

% \clearpage{\pagestyle{empty}\cleardoublepage}
% \include{glossario}


%\mainmatter
\acresetall
\chapter{Introdução}
\label{chap:intro}

\section{Enquadramento e Motivação}
\label{sec:amb} % CADA SECÇÃO DEVE TER UM LABEL
% CADA FIGURA DEVE TER UM LABEL
% CADA TABELA DEVE TER UM LABEL

Os acrónimos devem ser definidos recorrendo ao pacote (\emph{package}) \texttt{acronym}, usando os comandos \texttt{\textbackslash acro}, \texttt{\textbackslash ac}, \texttt{\textbackslash acp}, etc. E.g., \emph{The subject of this report is network protocols, namely is studied for several aspects of performance.}

Este relatório foi feito no contexto da unidade curricular de projeto da \ac{UBI}. Foi na \ac{UBI} que desenvolvi todo o trabalho.
                                                                        
\section{Objetivos}
\label{sec:obj}

\section{Organização do Documento}
\label{sec:organ}
% !POR EXEMPLO!
De modo a refletir o trabalho que foi feito, este documento encontra-se estruturado da seguinte forma:
\begin{enumerate}
\item O primeiro capítulo -- \textbf{Introdução} -- apresenta o projeto, a motivação para a sua escolha, o enquadramento para o mesmo, os seus objetivos e a respetiva organização do documento.
\item O segundo capítulo -- \textbf{Tecnologias Utilizadas} -- descreve os conceitos mais importantes no âmbito deste projeto, bem como as tecnologias utilizadas durante do desenvolvimento da aplicação.
\item ...
\end{enumerate}

% 
\section{Algumas Dicas -- [RETIRAR DA VERSÃO FINAL]}
% ALGUMAS DICAS
Os relatórios de projeto são individuais e preparados em \LaTeX, seguindo o formato disponível na página da unidade curricular. Deve ser prestada especial atenção aos seguintes pontos:
\begin{enumerate}
  \item O relatório deve ter um capítulo Introdução e Conclusões e Trabalho Futuro (ou só Conclusões);
  \item A última secção do primeiro capítulo deve descrever suscintamente a organização do documento;
  \item O relatório pode ser escrito em Língua Portuguesa ou Inglesa;
  \item Todas as imagens ou tabelas devem ter legendas e ser referidas no texto (usando comando \texttt{\textbackslash ref\{\}}).
\end{enumerate}
\chapter{Evolução da Governança da Internet}
% OU \chapter{Trabalhos Relacionados}
% OU \chapter{Engenharia de Software}
% OU \chapter{Tecnologias e Ferramentas Utilizadas}
\label{chap:evolução-da-governança-da-internet}

\section{Introdução}
\label{chap2:sec:intro}
Cada capítulo \underline{intermédio} deve começar com uma breve introdução onde é explicado com um pouco mais de detalhe qual é o tema deste capítulo, e como é que se encontra organizado (i.e., o que é que cada secção seguinte discute). 

%\section{Citações e Referências Cruzadas -- [RETIRAR DA VERSÃO FINAL]}
%\label{chap2:sec:citacoes}

Para se referenciarem outras secções, usar \texttt{\textbackslash{}ref\{label\}}, e.g., para citar a secção da Introdução deste capítulo, usar \texttt{\textbackslash{}ref\{chap2:sec:intro\}}. O resultado é: a secção~\ref{chap2:sec:intro} contém a introdução deste capítulo.

Para se citarem fontes bibliográficas, \underline{colocar a entrada certa} no ficheiro \texttt{bibiografia.bib} e usar o comando \texttt{\textbackslash{}cite\{label-da-referencia\}}, ligando o comando com a palavra que o antecede com um til. Por exemplo, para citar a referência eletrónica \emph{The Not So Short Introduction to \LaTeX{}}~\cite{short}, deve incluir-se o trecho seguinte no ficheiro \texttt{bibiografia.bib} e usar \texttt{\textbackslash{}cite\{\underline{short}\}} para a citação (citação incluída nesta mesma frase):
%
\small
\begin{verbatim}
@MISC{short,
  author = {Tobias Oetiker and Hubert Partl and 
            Irene Hyna and Elisabeth Schlegl},
  title  = "{The Not So Short Introduction to \LaTeX{}}",
  year   = 2018,
  note   = {[Online] \url{https://tobi.oetiker.ch/lshort/lshort.pdf}. 
            Último acesso a 12 de Março de 2019}
}
\end{verbatim}
\normalsize


\section{Secções Intermédias}
\label{chap2:sec:...}

\section{Conclusões}
\label{chap2:sec:concs}
Cada capítulo \underline{intermédio} deve referir o que demais importante se conclui desta parte do trabalho, de modo a fornecer a motivação para o capítulo ou passos seguintes.
\chapter{A Governança na atualidade}
% OU \chapter{Trabalhos Relacionados}
% OU \chapter{Engenharia de Software}
% OU \chapter{Tecnologias e Ferramentas Utilizadas}
\label{chap:tecno-ferra}

\section{Introdução}
\label{chap3:sec:intro}
Cada capítulo \underline{intermédio} deve começar com uma breve introdução onde é explicado com um pouco mais de detalhe qual é o tema deste capítulo, e como é que se encontra organizado (i.e., o que é que cada secção seguinte discute).

\section{Secções Intermédias}
\label{chap3:sec:...}

A tabela~\ref{tab:exemplo} serve apenas o propósito da exemplificação de como se fazem tabelas em \LaTeX.
%
\begin{table}
\centering
\begin{tabular}{|c|lr|}
\hline
\textbf{campo 1} & \textbf{campo 2} & \textbf{campo 3}\\
\hline
\hline
14 & 15 & 16 \\
\hline	
13 & 13 & 13 \\
\hline
\end{tabular}
\caption{Esta é uma tabela de exemplo.}
\label{tab:exemplo}
\end{table}

\section{Conclusões}
\label{chap3:sec:concs}
Cada capítulo \underline{intermédio} deve referir o que demais importante se conclui desta parte do trabalho, de modo a fornecer a motivação para o capítulo ou passos seguintes.
\chapter{Organização da Internet}
% Os titulos dados aos capítulos são meros exemplos. Cada relatório deve adequar-se ao projeto desenvolvido.
\label{chap:organização-da-internet}

\section{Introdução}
\label{chap4:sec:intro}
Cada capítulo \underline{intermédio} deve começar com uma breve introdução onde é explicado com um pouco mais de detalhe qual é o tema deste capítulo, e como é que se encontra organizado (i.e., o que é que cada secção seguinte discute).

\section{Secções Intermédias}
\label{chap4:sec:...}

O trecho de código seguinte mostra a função \texttt{main()} e o seu funcionamento:
\begin{lstlisting}[caption=Trecho de código usado no projeto.]
#include <stdio.h>

int main(){
  int i = 0;
  for(i = 0; i < 100; i++)
    printf("%d\n",i);
}
\end{lstlisting}


Se quiser definir a distribuição de Pareto, posso colocar a fórmula \emph{inline}, da seguinte forma $P(x)=\frac{x^{1/\Lambda}_{i}}{2}$, ou numa linha em separada, como se mostra a seguir:
$$ y^2 = \sum_{x=0}^{20}( x^3 - 2x + 3).$$

Outra maneira, mas numerada, é usar o ambiente \texttt{equation}, como se mostra na (\ref{eq:eq1}):
\begin{equation}
 y^2 = \sum_{x=0}^{20}( x^3 - 2x + 3).
 \label{eq:eq1}
\end{equation}

\begin{align}
 2+2+2+2+2+2+2+2+2+2+y^2 = & \sum_{x=0}^{20}( x^3 - 2x + 3);\\
                         = & x^4 -2.
 \label{eq:eq2}
\end{align}


\section{Conclusões}
\label{chap4:sec:concs}
Cada capítulo \underline{intermédio} deve referir o que demais importante se conclui desta parte do trabalho, de modo a fornecer a motivação para o capítulo ou passos seguintes.
\chapter{O caso Português}
\label{chap:caso-português}

\section{Introdução}
\label{chap5:sec:intro}


\include{capitulo-6}
\chapter*{Conclusão}
\addcontentsline{toc}{chapter}{Conclusão}
\label{chap:conc-trab-futuro}

Esta secção contém a resposta à questão: \\
\emph{Quais foram as conclusões princípais a que o(a) aluno(a) chegou no fim deste trabalho?}

Esta secção responde a questões como:\\
\emph{O que é que ficou por fazer, e porque?}\\
\emph{O que é que seria interessante fazer, mas não foi feito por não ser exatamente o objetivo deste trabalho?}\\
\emph{Em que outros casos ou situações ou cenários -- que não foram estudados no contexto deste projeto por não ser seu objetivo -- é que o trabalho aqui descrito pode ter aplicações interessantes e porque?}

% SE EXISTIREM APENDICES, DESCOMENTAR O QUE ESTÁ EM BAIXO
% \appendix
% \include{apendice1}
% \clearpage{\pagestyle{empty}\cleardoublepage}
% \include{continuacao}
% \clearpage{\pagestyle{empty}\cleardoublepage}
% \include{apendice2}
% \clearpage{\pagestyle{empty}\cleardoublepage}
% \include{apendice3}
% \clearpage{\pagestyle{empty}\cleardoublepage}

%\backmatter
\bibliographystyle{unsrt}
\bibliography{bibliografia}

\end{document}